% content-compact.tex --- leanblueprint for the Peter-Weyl theorem for
% compact (Hausdorff) topological groups, following T. Tao,
%   https://terrytao.wordpress.com/2011/01/23/the-peter-weyl-theorem-and-non-abelian-fourier-analysis-on-compact-groups/
%
% DESIGN: As detailed as possible, building from the ground up. Items are
% tagged for Mathlib readiness:
%   [ML]  = already in Mathlib (use \mathlibok and \leanok where applicable).
%   [~ML] = a trivial wrapper around something already in Mathlib.
%   [NEW] = not yet in Mathlib; needs local proof.
%
% Compared with the finite-group blueprint, the compact case requires
% substantial new infrastructure:
%   * Mathlib's `Representation` is purely algebraic; there is no built-in
%     "continuous unitary representation of a topological group on a Hilbert
%     space". We define one (UnitaryRep).
%   * Mathlib has `IsCompactOperator` (predicate) but **no spectral theorem
%     for compact self-adjoint operators in infinite dimensions**. We will
%     state the spectral theorem as [NEW]; an actual formalisation must
%     either prove it or use a workaround (e.g. via the cyclic-vector
%     reduction).
%   * Convolution on a compact group is in `Mathlib.Analysis.LConvolution`
%     and `Mathlib.Analysis.Convolution`, but the specific facts we need
%     (convolution by a continuous function is a compact operator on L^2,
%     etc.) are mostly [NEW].
%
% Key Mathlib files used:
%   Mathlib/MeasureTheory/Group/Measure.lean          IsHaarMeasure, IsMulLeftInvariant
%   Mathlib/MeasureTheory/Measure/Haar/{Basic, Unique}.lean
%   Mathlib/MeasureTheory/Function/L2Space.lean       L^2(G), inner product
%   Mathlib/Analysis/InnerProductSpace/{Defs, Basic, l2Space}.lean
%   Mathlib/Analysis/Normed/Operator/Compact.lean     IsCompactOperator
%   Mathlib/Topology/ContinuousMap/StoneWeierstrass.lean
%   Mathlib/Analysis/Convolution.lean, Mathlib/Analysis/LConvolution.lean
%
% Standing setup throughout:
%   G : Type*, [Group G] [TopologicalSpace G] [TopologicalGroup G]
%       [CompactSpace G] [T2Space G] [MeasurableSpace G] [BorelSpace G]
%   μ : Measure G with [μ.IsHaarMeasure] and probability normalisation
%       (μ univ = 1).
%   k : RCLike (= ℝ or ℂ); throughout we work over ℂ.

\chapter{Setup, Haar Measure, and Hilbert-Space Framework}
\label{chap:c-setup}

\section{Standing hypotheses}

\begin{definition}[Standing hypotheses for compact Peter--Weyl]
  \label{def:c-standing}
  \mathlibok
  Throughout the blueprint, $G$ denotes a Hausdorff compact topological
  group, written multiplicatively, with identity $1_G$. In Lean we use
  \begin{center}
    \texttt{[Group G] [TopologicalSpace G] [IsTopologicalGroup G]
      [CompactSpace G] [T2Space G] [MeasurableSpace G] [BorelSpace G]}.
  \end{center}
  The scalar field is $k = \mathbb C$ (more precisely
  \texttt{[RCLike k]} so the same blueprint covers the real case after the
  obvious complexification step).  The Haar measure is denoted $\mu$ and
  is normalised to be a probability measure ($\mu(G) = 1$).
\end{definition}

\section{Haar measure on a compact group}

\begin{theorem}[Existence and uniqueness of Haar measure]
  \label{thm:c-haar-exists}
  \lean{MeasureTheory.Measure.haarMeasure, MeasureTheory.Measure.IsHaarMeasure}
  \mathlibok
  \leanok
  \uses{def:c-standing}
  Every locally compact Hausdorff topological group $G$ admits a left-invariant
  regular Borel measure that is finite on compact sets and positive on
  non-empty open sets, unique up to a positive scalar.  In Lean this is
  the constant \texttt{MeasureTheory.Measure.haarMeasure} (the file
  \texttt{Mathlib.MeasureTheory.Measure.Haar.Basic}); the typeclass
  \texttt{IsHaarMeasure $\mu$} packages all the invariance and regularity
  conditions; uniqueness is in
  \texttt{Mathlib.MeasureTheory.Measure.Haar.Unique} via
  \texttt{haarScalarFactor} and
  \texttt{isMulLeftInvariant\_eq\_smul\_of\_regular}.
\end{theorem}

\begin{proof}\leanok\uses{def:c-standing}
  Direct citation of the existence and uniqueness theorems in Mathlib's
  Haar files.
\end{proof}

\begin{corollary}[Haar measure is finite for compact $G$]
  \label{cor:c-haar-finite}
  \lean{MeasureTheory.IsFiniteMeasure}
  \mathlibok
  \leanok
  \uses{thm:c-haar-exists}
  Since $G$ is compact and $\mu$ is finite on compact sets, $\mu(G) < \infty$.
  We may rescale so that $\mu(G) = 1$ (probability normalisation), and we
  do so throughout.
\end{corollary}

\begin{proof}\leanok
  \uses{thm:c-haar-exists}
  \texttt{IsFiniteMeasureOnCompacts} is part of the \texttt{IsHaarMeasure}
  bundle. Rescale via
  \texttt{MeasureTheory.Measure.haarScalarFactor}.
\end{proof}

\begin{theorem}[Right-invariance of Haar measure on compact $G$]
  \label{thm:c-haar-right-invariant}
  \lean{MeasureTheory.Measure.IsMulRightInvariant}
  \mathlibok
  \leanok
  \uses{thm:c-haar-exists, cor:c-haar-finite}
  On a compact (Hausdorff) topological group, every left Haar measure is
  also right-invariant: $\mu(Sg) = \mu(S)$ for every measurable $S$ and
  $g \in G$. (Equivalently, compact groups are unimodular.)
\end{theorem}

\begin{proof}\leanok
  \uses{thm:c-haar-exists}
  Mathlib's \texttt{IsHaarMeasure.IsMulRightInvariant} (or the
  \texttt{IsHaarMeasure.unimodular}-style lemma) handles this. The
  classical proof: the modular function $\Delta : G \to (0, \infty)$ is a
  continuous group homomorphism; on a compact group its image is a
  compact subgroup of $(0, \infty)$ which is just $\{1\}$, so
  $\Delta \equiv 1$.
\end{proof}

\begin{theorem}[Inversion-invariance of Haar measure]
  \label{thm:c-haar-inv}
  \lean{MeasureTheory.Measure.inv}
  \mathlibok
  \leanok
  \uses{thm:c-haar-right-invariant}
  Pushforward of $\mu$ under $g \mapsto g^{-1}$ equals $\mu$.
\end{theorem}

\begin{proof}\leanok
  \uses{thm:c-haar-right-invariant}
  \texttt{Measure.inv\_eq\_self\_of\_isHaarMeasure} (or the corresponding
  Mathlib name for unimodular groups).
\end{proof}

\section{The Hilbert space $L^2(G)$}

\begin{definition}[$L^2(G)$]
  \label{def:c-L2}
  \lean{MeasureTheory.Lp}
  \mathlibok
  \uses{def:c-standing, thm:c-haar-exists}
  $L^2(G) := \texttt{Lp $\mathbb C$ 2 $\mu$}$, the Lebesgue $L^2$ space of
  $\mathbb C$-valued square-integrable functions on $G$ modulo a.e.\
  equality, equipped with the inner product
  $\langle f, h \rangle := \int_G f(x)\overline{h(x)}\,d\mu(x)$.
  In Lean this is from \texttt{Mathlib.MeasureTheory.Function.L2Space}; the
  typeclass instance \texttt{L2.innerProductSpace} provides the inner
  product, and \texttt{MeasureTheory.Lp.completeSpace} provides
  completeness.
\end{definition}

\begin{theorem}[$C(G)$ is dense in $L^2(G)$]
  \label{thm:c-CG-dense}
  \lean{MeasureTheory.Lp.continuousMap_dense}
  \mathlibok
  \leanok
  \uses{def:c-L2}
  Continuous functions $C(G) \hookrightarrow L^2(G)$ form a dense subspace.
\end{theorem}

\begin{proof}\leanok\uses{def:c-L2}
  Mathlib's general
  \texttt{MeasureTheory.Lp.continuousMap\_dense} (via mollification /
  outer regularity); valid for any finite regular Borel measure on a
  Polish space, hence in particular for compact Hausdorff $G$.
\end{proof}

\begin{theorem}[$C(G) \subseteq L^\infty(G) \subseteq L^2(G) \subseteq L^1(G)$]
  \label{thm:c-Lp-inclusions}
  \lean{MeasureTheory.MemLp.mono_exponent_of_isFiniteMeasure}
  \mathlibok
  \leanok
  \uses{cor:c-haar-finite, def:c-L2}
  Since $\mu$ is finite, $L^p$ inclusions go the natural way:
  $\|f\|_{L^p} \leq \mu(G)^{1/p - 1/q} \|f\|_{L^q}$ for $p \leq q$.
  Continuous functions on a compact $G$ are bounded so are in
  $L^\infty$.
\end{theorem}

\begin{proof}\leanok\uses{cor:c-haar-finite}
  Mathlib's \texttt{MeasureTheory.MemLp.mono\_exponent\_of\_isFiniteMeasure}.
\end{proof}


\chapter{Continuous Unitary Representations}
\label{chap:c-rep}

Mathlib's \texttt{Representation k G V} is purely algebraic: there is no
continuity assumption and no Hilbert-space structure. We must add a layer
of definitions for the Peter--Weyl theorem.  All items in this chapter are
\texttt{[NEW]} unless flagged otherwise.

\begin{definition}[Continuous unitary representation]
  \label{def:c-unitary-rep}
  \lean{UnitaryRep}
  % [NEW] Local definition; not in Mathlib.
  \uses{def:c-standing}
  A \emph{continuous unitary representation} of $G$ on a (complex) Hilbert
  space $H$ is a group homomorphism $\rho : G \to U(H)$ to the unitary
  group of $H$, that is continuous when $U(H)$ is given the strong
  operator topology.  In Lean we encode this as a structure
  \texttt{UnitaryRep $\mathbb C$ G H} bundling
  \texttt{toMonoidHom : G $\to$* (H $\to$L[$\mathbb C$] H)}, an
  \texttt{isUnitary} field, and a \texttt{continuous} field with respect
  to the strong operator topology on
  \texttt{H $\to$L[$\mathbb C$] H}.
\end{definition}

\begin{remark}
  Equivalently, by a standard argument, strong operator continuity of
  $\rho$ is implied by either of: (a) continuity of
  $g \mapsto \rho(g) v$ for every $v \in H$; (b) continuity of
  $(g, v) \mapsto \rho(g) v$. We will use (a) as our default working
  definition; see Lemma~\ref{lem:c-strong-vs-weak-cts}.
\end{remark}

\begin{lemma}[Strong vs weak operator continuity]
  \label{lem:c-strong-vs-weak-cts}
  \lean{UnitaryRep.continuous_iff_continuous_apply}
  % [NEW]
  \uses{def:c-unitary-rep}
  For a homomorphism $\rho : G \to U(H)$, the following are equivalent:
  \begin{enumerate}
    \item $\rho$ is continuous in the strong operator topology;
    \item for every $v \in H$ the orbit map $g \mapsto \rho(g) v$ is
      continuous;
    \item the action map $G \times H \to H$, $(g, v) \mapsto \rho(g) v$
      is jointly continuous.
  \end{enumerate}
\end{lemma}

\begin{proof}
  \uses{def:c-unitary-rep}
  Standard. (1) $\Leftrightarrow$ (2) is the definition of the strong
  operator topology. (2) $\Rightarrow$ (3) uses uniform boundedness of
  $\|\rho(g)\| = 1$ together with the $\varepsilon/3$ argument.
  (3) $\Rightarrow$ (2) is obvious.
\end{proof}

\begin{definition}[The left regular representation $\lambda$]
  \label{def:c-regular-rep}
  \lean{UnitaryRep.leftRegular}
  % [NEW]
  \uses{def:c-L2, def:c-unitary-rep}
  The \emph{left regular representation} $\lambda : G \to U(L^2(G))$ is
  defined by $(\lambda(g) f)(x) := f(g^{-1} x)$ on a continuous
  representative, extended by continuity to all of $L^2(G)$.
\end{definition}

\begin{theorem}[$\lambda$ is a continuous unitary representation]
  \label{thm:c-regular-unitary-cts}
  \lean{UnitaryRep.leftRegular_isUnitaryRep}
  % [NEW]
  \uses{def:c-regular-rep, thm:c-haar-exists, thm:c-haar-right-invariant}
  $\lambda(g)$ is unitary on $L^2(G)$ for each $g$ (by left-invariance of
  $\mu$). The map $g \mapsto \lambda(g)$ is strongly continuous (using
  density of $C(G)$ in $L^2(G)$ together with uniform continuity of any
  continuous function on a compact group).
\end{theorem}

\begin{proof}
  \uses{thm:c-haar-exists, thm:c-CG-dense}
  Unitarity: change of variables under left translation preserves $\mu$
  by Theorem~\ref{thm:c-haar-exists}.
  Strong continuity: for $f \in C(G)$, the orbit map
  $g \mapsto \lambda(g) f$ is continuous in $L^2$ by uniform continuity
  of $f$ on the compact group $G$ (Heine--Cantor); for general
  $f \in L^2$, use density of $C(G)$ from Theorem~\ref{thm:c-CG-dense}
  and the uniform bound $\|\lambda(g)\| = 1$ via an
  $\varepsilon/3$-argument.
\end{proof}

\begin{definition}[Right regular representation]
  \label{def:c-right-reg}
  \lean{UnitaryRep.rightRegular}
  % [NEW]
  \uses{def:c-L2, def:c-unitary-rep}
  $(\rho(g) f)(x) := f(xg)$. Same proofs as for $\lambda$ show $\rho$ is
  a continuous unitary representation. The two representations
  \emph{commute}: $\lambda(g)\rho(h) = \rho(h)\lambda(g)$ for all
  $g, h \in G$.
\end{definition}

\begin{definition}[Subrepresentation]
  \label{def:c-subrep}
  \lean{UnitaryRep.Subrep}
  % [NEW]
  \uses{def:c-unitary-rep}
  A \emph{subrepresentation} of a unitary representation $(\rho, H)$ is a
  closed $\rho$-invariant subspace $V \leq H$. Subrepresentations are
  closed under taking orthogonal complements (since unitary operators
  preserve orthogonality and closed subspaces).
\end{definition}

\begin{lemma}[Orthogonal complement of a subrepresentation]
  \label{lem:c-subrep-orthcomp}
  \lean{UnitaryRep.Subrep.orthogonalComplement}
  % [NEW]
  \uses{def:c-subrep}
  If $V$ is a closed $\rho$-invariant subspace, then $V^\perp$ is also
  $\rho$-invariant. Hence $H = V \oplus V^\perp$ as a direct sum of
  subrepresentations.
\end{lemma}

\begin{proof}
  \uses{def:c-subrep}
  $\rho(g)$ is unitary, so it preserves the inner product:
  if $w \in V^\perp$ and $v \in V$, then
  $\langle \rho(g)w, v \rangle = \langle w, \rho(g^{-1})v \rangle = 0$
  since $\rho(g^{-1})v \in V$. Thus $\rho(g)w \in V^\perp$.
\end{proof}

\begin{definition}[Irreducible representation]
  \label{def:c-irrep}
  \lean{UnitaryRep.IsIrreducible}
  % [NEW]
  \uses{def:c-subrep}
  A continuous unitary representation $(\rho, H)$ is \emph{irreducible}
  if $H \neq 0$ and the only closed $\rho$-invariant subspaces are
  $\{0\}$ and $H$ itself.
\end{definition}

\begin{definition}[Intertwiner]
  \label{def:c-intertwiner}
  \lean{UnitaryRep.Intertwiner}
  % [NEW]
  \uses{def:c-unitary-rep}
  An \emph{intertwiner} between two continuous unitary representations
  $(\rho, H)$ and $(\sigma, K)$ is a bounded linear map $T : H \to K$
  with $T \rho(g) = \sigma(g) T$ for all $g$. We write
  $\mathrm{Hom}_G(\rho, \sigma)$ for the space of such intertwiners.
\end{definition}

\begin{definition}[Equivalence of representations]
  \label{def:c-equiv-rep}
  \lean{UnitaryRep.Equiv}
  % [NEW]
  \uses{def:c-intertwiner}
  $(\rho, H) \cong (\sigma, K)$ if there exists a unitary intertwiner
  $H \to K$.
\end{definition}


\chapter{Schur's Lemma in the Compact Setting}
\label{chap:c-schur}

\section{Schur's lemma for irreducible unitary representations}

\begin{theorem}[Schur's lemma, scalar form]
  \label{thm:c-schur-scalar}
  \lean{UnitaryRep.IsIrreducible.intertwiner_self_eq_scalar}
  % [NEW]
  \uses{def:c-irrep, def:c-intertwiner}
  Let $(\rho, H)$ be an irreducible continuous unitary representation
  of $G$. Then every bounded $G$-equivariant operator $T : H \to H$ is
  a scalar: $T = \lambda \cdot \mathrm{Id}_H$ for some
  $\lambda \in \mathbb C$.
\end{theorem}

\begin{proof}
  \uses{def:c-irrep, def:c-intertwiner, lem:c-subrep-orthcomp}
  Decompose $T = A + iB$ with $A := (T+T^*)/2$ and $B := (T-T^*)/(2i)$
  self-adjoint. Each of $A, B$ commutes with $\rho$. Apply the spectral
  theorem to $A$: every eigenspace $\ker(A - \lambda I)$ is closed and
  $\rho$-invariant, so by irreducibility it is $\{0\}$ or $H$. Hence the
  spectrum of $A$ is a single point and $A$ is a scalar; similarly for
  $B$.

  (For the formalisation: this proof needs the bounded-self-adjoint
  spectral theorem. A cleaner Mathlib-friendly route is to argue via
  cyclic vectors using only the existence of a non-trivial reducing
  projection from approximate eigenspaces, but we follow the standard
  spectral-theorem approach.)
\end{proof}

\begin{corollary}[Schur's lemma, distinct irreducibles]
  \label{cor:c-schur-distinct}
  \lean{UnitaryRep.IsIrreducible.subsingleton_intertwiner_of_not_equiv}
  % [NEW]
  \uses{thm:c-schur-scalar, def:c-equiv-rep}
  If $(\rho, H)$ and $(\sigma, K)$ are irreducible and not isomorphic,
  then $\mathrm{Hom}_G(\rho, \sigma) = 0$.
\end{corollary}

\begin{proof}
  \uses{thm:c-schur-scalar}
  Let $T \in \mathrm{Hom}_G(\rho, \sigma)$. Then $\overline{\ker T}$
  is a closed $\rho$-invariant subspace of $H$, so it is $0$ or $H$.
  Similarly $\overline{T(H)}$ is $0$ or $K$. If $T \neq 0$, then
  $\ker T = 0$ and $\overline{T(H)} = K$. Polar decomposition $T = U|T|$
  expresses $T$ as a unitary $U : H \to K$ times a positive operator
  $|T| := \sqrt{T^*T}$ on $H$. By $G$-equivariance of $T$, both $U$
  and $|T|$ are $G$-equivariant. The operator $|T| \in \mathrm{Hom}_G(\rho, \rho)$
  is then a positive scalar by Theorem~\ref{thm:c-schur-scalar}, and $U$ is a
  unitary intertwiner $\rho \to \sigma$, contradicting non-equivalence.
\end{proof}


\chapter{Convolution and Compact Self-adjoint Operators}
\label{chap:c-convolution}

This is the analytic core of Tao's proof.  Mathlib has convolution
defined (\texttt{Mathlib.Analysis.LConvolution}, additive notation;
\texttt{Mathlib.Analysis.Convolution}), but most of the specific facts
we need below are \texttt{[NEW]} for compact non-abelian groups.

\section{Convolution on a compact group}

\begin{definition}[Convolution on a compact group]
  \label{def:c-conv}
  \lean{MeasureTheory.mconvolution}
  \mathlibok
  \uses{def:c-standing, thm:c-haar-exists}
  For $f, g : G \to \mathbb C$ Borel-measurable, the (multiplicative)
  convolution is
  \[
     (f * g)(x) := \int_G f(y) g(y^{-1}x) \, d\mu(y).
  \]
  Mathlib has \texttt{MeasureTheory.mconvolution} (in
  \texttt{Mathlib.Analysis.Convolution}, set up via a continuous bilinear
  form $L$) for general topological groups.
\end{definition}

\begin{theorem}[Young's inequality on a compact group]
  \label{thm:c-young}
  \lean{MeasureTheory.MemLp.convolution}
  \mathlibok
  \leanok
  \uses{def:c-conv, def:c-L2, thm:c-Lp-inclusions}
  For $\mu(G) = 1$:
  $\|f * g\|_{L^2} \leq \|f\|_{L^1}\|g\|_{L^2}$ and
  $\|f * g\|_{L^\infty} \leq \|f\|_{L^2}\|g\|_{L^2}$. Hence
  $L^2(G) * L^2(G) \subseteq L^\infty(G) \cap L^2(G) \subseteq C(G)$
  (the last when $f \in L^2$ and $g \in C(G)$, or via uniform
  approximation).
\end{theorem}

\begin{proof}
  \uses{def:c-conv}
  Direct computation; in Mathlib this follows from versions of Young's
  inequality in \texttt{MeasureTheory.MemLp.convolution} after specialising
  to compact groups.
\end{proof}

\begin{theorem}[Convolution by an $L^2$ kernel is compact on $L^2(G)$]
  \label{thm:c-conv-compact}
  \lean{MeasureTheory.convolution_isCompactOperator}
  % [NEW] Standard, but not packaged for general compact groups in Mathlib.
  \uses{thm:c-young, def:c-L2}
  Let $\phi \in L^2(G)$. Then the operator
  $T_\phi : L^2(G) \to L^2(G)$, $f \mapsto f * \phi$, is a Hilbert--Schmidt
  operator with kernel $K(x, y) = \phi(y^{-1}x) \in L^2(G \times G)$. In
  particular $T_\phi$ is a compact operator (in the sense of Mathlib's
  \texttt{IsCompactOperator}).
\end{theorem}

\begin{proof}
  \uses{thm:c-young}
  The operator $T_\phi$ is given by an integral kernel
  $K(x, y) = \phi(y^{-1}x)$. Compute its Hilbert--Schmidt norm:
  \[
     \int\int |K(x,y)|^2 d\mu\,d\mu
     = \int\int |\phi(y^{-1}x)|^2 d\mu(y)\,d\mu(x)
     = \mu(G)\int |\phi|^2 = \|\phi\|_{L^2}^2 < \infty,
  \]
  using right-invariance (Theorem~\ref{thm:c-haar-right-invariant}) and
  inversion-invariance of $\mu$. Hilbert--Schmidt operators are compact:
  this requires the (currently unformalised in Mathlib for the general
  Hilbert-Schmidt case) fact that integral operators with $L^2$ kernels
  are compact. Equivalently: the Hilbert-Schmidt norm bounds the operator
  norm by the trace-class lemma.
\end{proof}

\begin{theorem}[Self-adjointness of convolution]
  \label{thm:c-conv-selfadjoint}
  \lean{MeasureTheory.convolution_isSelfAdjoint}
  % [NEW]
  \uses{thm:c-conv-compact, thm:c-haar-inv}
  Let $\phi \in L^2(G)$ satisfy $\phi(g) = \overline{\phi(g^{-1})}$ for
  $\mu$-a.e.\ $g$. Then the operator
  $T_\phi : f \mapsto f * \phi$ is self-adjoint on $L^2(G)$.
\end{theorem}

\begin{proof}
  \uses{thm:c-haar-inv}
  Compute $\langle T_\phi f, h \rangle$ and
  $\langle f, T_\phi h \rangle$ via Fubini. The condition
  $\phi(g) = \overline{\phi(g^{-1})}$ exactly switches the two integrands
  after a substitution $g \mapsto g^{-1}$ (using inversion invariance,
  Theorem~\ref{thm:c-haar-inv}).
\end{proof}

\begin{theorem}[Convolution commutes with right translation]
  \label{thm:c-conv-comm-right}
  \lean{MeasureTheory.convolution_commutes_rightRegular}
  % [NEW]
  \uses{def:c-conv, def:c-right-reg}
  For every $\phi \in L^2(G)$ and $g \in G$,
  $T_\phi \circ \rho(g) = \rho(g) \circ T_\phi$.
\end{theorem}

\begin{proof}
  \uses{def:c-conv, def:c-right-reg}
  Direct change of variables: for $f \in L^2(G)$,
  $T_\phi(\rho(g)f)(x) = \int f(yg) \phi(y^{-1}x) d\mu(y)
   = \int f(z) \phi(z^{-1} gx) d\mu(z)$
  (substituting $z = yg$, using Theorem~\ref{thm:c-haar-right-invariant}),
  which is $(T_\phi f)(gx) = \rho(g)(T_\phi f)(x)$.
\end{proof}

\section{Spectral theorem for compact self-adjoint operators}

\begin{theorem}[Spectral theorem for compact self-adjoint operators]
  \label{thm:c-spectral}
  \lean{ContinuousLinearMap.IsCompactOperator.IsSelfAdjoint.eigenspaceDecomposition}
  % [NEW] Mathlib has the *finite-dimensional* spectral theorem
  %       (Mathlib.Analysis.InnerProductSpace.Spectrum) but as of the
  %       knowledge cutoff does not have the spectral theorem for compact
  %       self-adjoint operators in the Hilbert-space setting.
  \uses{thm:c-conv-compact, thm:c-conv-selfadjoint}
  Let $T$ be a non-zero compact self-adjoint operator on a complex
  Hilbert space $H$. Then:
  \begin{enumerate}
    \item the spectrum $\sigma(T) \setminus \{0\}$ consists of countably
      many real eigenvalues $\{\lambda_n\}$ with $\lambda_n \to 0$;
    \item each non-zero eigenvalue has finite multiplicity, i.e.\ the
      eigenspace $H_{\lambda_n} := \ker(T - \lambda_n I)$ is
      finite-dimensional;
    \item we have an orthogonal Hilbert-sum decomposition
      $H = \ker T \oplus \widehat{\bigoplus}_n H_{\lambda_n}$.
  \end{enumerate}
  Equivalently: there exists a (countable) orthonormal basis of $H$
  consisting of eigenvectors of $T$.
\end{theorem}

\begin{proof}
  \uses{thm:c-conv-compact}
  Standard. Sketch: the operator norm $\|T\|$ is attained as
  $\sup\{|\langle Tx, x \rangle|\} = \pm\lambda_1$, an eigenvalue, with
  finite-dimensional eigenspace $H_{\lambda_1}$; pass to
  $H_{\lambda_1}^\perp$ which is $T$-invariant, and iterate. The
  iterating eigenvalues form a sequence going to zero by compactness
  (any bounded set of eigenvectors with eigenvalues bounded below has
  a convergent subsequence, but eigenvectors of distinct eigenvalues
  are orthogonal, so a uniformly lower-bounded sequence cannot have
  a convergent subsequence). The closure of the span of all eigenvectors
  with non-zero eigenvalue equals $(\ker T)^\perp$.

  This is a non-trivial result of functional analysis that, to the
  author's knowledge, is not currently in Mathlib (as of knowledge
  cutoff Jan 2026; please verify against current Mathlib state). The
  finite-dimensional version
  \texttt{LinearMap.IsSymmetric.diagonalization} is in
  \texttt{Mathlib.Analysis.InnerProductSpace.Spectrum}; extending to
  the compact-operator case is itself a substantial formalisation
  project.
\end{proof}


\chapter{Finite-Dimensionality of Eigenspaces and Existence of Irreducibles}
\label{chap:c-finite-dim}

\section{Eigenspaces of $G$-equivariant compact operators}

\begin{theorem}[Eigenspaces of equivariant compact operators are
  $G$-invariant]
  \label{thm:c-eigenspace-invariant}
  \lean{UnitaryRep.eigenspace_isInvariant}
  % [NEW]
  \uses{thm:c-conv-comm-right, thm:c-spectral}
  Let $\phi \in L^2(G)$ and $T_\phi$ the convolution operator on
  $L^2(G)$. For each $\lambda \in \mathbb C$, the eigenspace
  $H_\lambda := \ker(T_\phi - \lambda I)$ is closed and right-translation-invariant
  (i.e.\ a closed subrepresentation of the right regular representation).
\end{theorem}

\begin{proof}
  \uses{thm:c-conv-comm-right}
  $T_\phi$ commutes with $\rho(g)$ (Theorem~\ref{thm:c-conv-comm-right}),
  so if $T_\phi v = \lambda v$ then $T_\phi(\rho(g)v) = \rho(g)T_\phi v
  = \lambda \rho(g) v$, i.e.\ $\rho(g)v \in H_\lambda$.
\end{proof}

\begin{corollary}[Every irreducible subrepresentation of $L^2(G)$ is
  finite-dimensional]
  \label{cor:c-irr-finite-dim}
  \lean{UnitaryRep.subrepresentation_finiteDimensional_of_irreducible}
  % [NEW]
  \uses{thm:c-eigenspace-invariant, thm:c-spectral, thm:c-conv-selfadjoint}
  Let $V \leq L^2(G)$ be a non-zero irreducible right-translation-invariant
  closed subspace. Then $\dim V < \infty$.
\end{corollary}

\begin{proof}
  \uses{thm:c-eigenspace-invariant, thm:c-spectral, thm:c-conv-selfadjoint}
  Let $0 \neq v \in V$. Pick $\phi \in C(G)$ self-adjoint
  ($\phi(g) = \overline{\phi(g^{-1})}$) such that $T_\phi v \neq 0$
  (such $\phi$ exists: e.g.\ a small bump function around $1_G$
  approximates the identity, see Lemma~\ref{lem:c-approx-identity}).
  By Theorem~\ref{thm:c-conv-compact} and
  Theorem~\ref{thm:c-conv-selfadjoint}, $T_\phi$ is compact and
  self-adjoint, and by the spectral theorem
  (Theorem~\ref{thm:c-spectral}) it has at least one non-zero eigenvalue
  $\lambda$, with $H_\lambda$ finite-dimensional.
  By Theorem~\ref{thm:c-eigenspace-invariant} the eigenspace
  $H_\lambda$ is right-translation-invariant, so
  $V \cap H_\lambda$ is a non-zero (since $T_\phi v \neq 0$ for some
  eigenfunction's spectral component lies in $V$) right-invariant
  subspace of $V$.  By irreducibility of $V$, $V \cap H_\lambda = V$,
  so $V \subseteq H_\lambda$, which is finite-dimensional.
\end{proof}

\begin{lemma}[Approximate identity from continuous bumps]
  \label{lem:c-approx-identity}
  \lean{ContinuousMap.bumpFunction_approxIdentity}
  % [NEW]
  \uses{def:c-conv, thm:c-haar-exists}
  There exists a net $\{\phi_U\}_{U \in \mathcal{N}(1_G)}$ of continuous
  non-negative functions $\phi_U : G \to \mathbb R_{\geq 0}$, indexed by
  open neighbourhoods $U$ of $1_G$, satisfying:
  $\mathrm{supp}(\phi_U) \subseteq U$,
  $\int \phi_U \, d\mu = 1$,
  $\phi_U(g) = \phi_U(g^{-1})$, and
  $f * \phi_U \to f$ in $L^2$ as $U \to \{1_G\}$ for every
  $f \in L^2(G)$.
\end{lemma}

\begin{proof}
  \uses{thm:c-haar-exists}
  Use Urysohn's lemma to obtain continuous non-negative bump functions
  with support in $U$, normalise the integral, symmetrise via
  $\phi_U(g) := \tfrac12(\psi_U(g) + \psi_U(g^{-1}))$. Convergence
  $f * \phi_U \to f$ uses uniform continuity of $f \in C(G)$ on the
  compact group $G$ followed by density of $C(G)$ in $L^2$.
\end{proof}


\chapter{The Peter--Weyl Decomposition of $L^2(G)$}
\label{chap:c-peter-weyl-L2}

\section{Matrix coefficients}

\begin{definition}[Matrix coefficient]
  \label{def:c-matrix-coeff}
  \lean{UnitaryRep.matrixCoeff}
  % [NEW]
  \uses{def:c-unitary-rep}
  Given a finite-dimensional continuous unitary representation
  $(\rho, V)$ of $G$ and vectors $u, v \in V$, the function
  $\pi_{u,v}(g) := \langle \rho(g) u, v \rangle_V$
  is the \emph{matrix coefficient} associated to $(\rho, u, v)$.
  Matrix coefficients are continuous functions on $G$ and bounded
  (by $\|u\|\,\|v\|$). Hence $\pi_{u,v} \in C(G) \cap L^2(G)$.
\end{definition}

\begin{lemma}[Algebra structure on matrix coefficients]
  \label{lem:c-mc-algebra}
  \lean{UnitaryRep.matrixCoeff_algebra}
  % [NEW]
  \uses{def:c-matrix-coeff}
  The set of all matrix coefficients (over all finite-dimensional
  continuous unitary representations of $G$) is closed under: pointwise
  addition (taking direct sums); pointwise multiplication (taking
  tensor products); complex conjugation (taking the contragredient
  $\overline{\rho}$); and contains the constants (the trivial
  representation).
\end{lemma}

\begin{proof}
  \uses{def:c-matrix-coeff}
  Standard verifications: e.g.\ if
  $f_1 = \pi^{(1)}_{u_1, v_1}$ for $\rho_1$ on $V_1$ and
  $f_2 = \pi^{(2)}_{u_2, v_2}$ for $\rho_2$ on $V_2$, then
  $f_1 \cdot f_2 = \pi^{(\rho_1 \otimes \rho_2)}_{u_1 \otimes u_2,
                                                 v_1 \otimes v_2}$.
  Conjugation $\overline{f_1}$ is a matrix coefficient of the
  contragredient representation on $V_1^*$.
\end{proof}

\begin{definition}[The matrix-coefficient subalgebra $\mathcal A$]
  \label{def:c-mc-subalg}
  \lean{UnitaryRep.matrixCoeffSubalgebra}
  % [NEW]
  \uses{lem:c-mc-algebra}
  Let $\mathcal A := \mathrm{span}_{\mathbb C}
   \{\pi_{u,v} : (\rho, V) \text{ fin-dim cts unitary rep},\ u, v \in V\}
   \subseteq C(G)$.
  By Lemma~\ref{lem:c-mc-algebra}, $\mathcal A$ is a unital
  $*$-subalgebra of the (commutative under pointwise multiplication)
  $*$-algebra $C(G)$. Equivalently: a star-subalgebra of the Mathlib
  bundled type \texttt{StarSubalgebra $\mathbb C$ C(G, $\mathbb C$)}.
\end{definition}

\section{$\mathcal A$ separates points: existence of irreducibles}

This is the key step in proving Peter--Weyl. Tao calls the underlying
proposition ``every non-trivial element of $G$ acts non-trivially on some
finite-dimensional irreducible''.

\begin{theorem}[Existence of separating finite-dimensional irreducibles]
  \label{thm:c-irreducibles-separate}
  \lean{UnitaryRep.separatesPoints}
  % [NEW]
  \uses{thm:c-conv-compact, thm:c-spectral, cor:c-irr-finite-dim,
        thm:c-eigenspace-invariant, lem:c-approx-identity}
  For every $g \in G$ with $g \neq 1_G$, there exists a finite-dimensional
  continuous unitary irreducible representation $(\rho, V)$ of $G$ with
  $\rho(g) \neq \mathrm{Id}_V$.
\end{theorem}

\begin{proof}
  \uses{thm:c-conv-compact, thm:c-spectral, cor:c-irr-finite-dim,
        lem:c-approx-identity, def:c-regular-rep}
  Following Tao. Pick $g \neq 1_G$. By compactness and Hausdorffness,
  pick disjoint open neighbourhoods $U \ni 1_G$ and $V \ni g$, and a
  continuous symmetric non-negative bump $\phi$ supported in $U$ with
  $\int \phi = 1$.

  Then $T_\phi(\mathbf{1}_{V}) = \mathbf{1}_V * \phi$ is supported in
  $V \cdot U$ (which is disjoint from a smaller neighbourhood of
  $1_G$), so the function $T_\phi(\mathbf{1}_V) - \mathbf{1}_V$ is
  non-zero. Pick a non-zero eigenvalue $\lambda$ of $T_\phi$ such that
  $H_\lambda$ is not contained in the $g$-fixed subspace of $L^2(G)$
  (this is possible since the spans of all $H_\lambda$ are dense in
  $(\ker T_\phi)^\perp$ by Theorem~\ref{thm:c-spectral}, and $g$ acts
  non-trivially on $L^2(G)$ via the right regular representation
  $\rho(g)$, since $\rho(g) \mathbf{1}_V \neq \mathbf{1}_V$ when
  $V$ is small).

  By Theorem~\ref{thm:c-eigenspace-invariant} the eigenspace
  $H_\lambda$ is finite-dimensional and right-$G$-invariant.  Decompose
  $H_\lambda$ into irreducible subrepresentations (possible since
  $H_\lambda$ is a finite-dimensional unitary representation, which is
  semisimple by Lemma~\ref{lem:c-subrep-orthcomp}). At least one
  irreducible component $V \subseteq H_\lambda$ has $\rho(g) \neq
  \mathrm{Id}_V$ (else $g$ acts trivially on $H_\lambda$, contradicting
  the choice).
\end{proof}

\begin{corollary}[$\mathcal A$ separates points]
  \label{cor:c-mc-sep-points}
  \lean{UnitaryRep.matrixCoeffSubalgebra_separatesPoints}
  % [NEW]
  \uses{thm:c-irreducibles-separate, def:c-mc-subalg}
  The matrix-coefficient $*$-subalgebra
  $\mathcal A \subseteq C(G, \mathbb C)$ separates points of $G$.
\end{corollary}

\begin{proof}
  \uses{thm:c-irreducibles-separate, def:c-mc-subalg}
  Let $g \neq h$ in $G$. Apply Theorem~\ref{thm:c-irreducibles-separate}
  to $gh^{-1}$ to obtain a finite-dimensional irreducible $(\rho, V)$
  with $\rho(gh^{-1}) \neq \mathrm{Id}_V$. Pick $u, v \in V$ with
  $\langle \rho(gh^{-1}) u, v \rangle \neq \langle u, v \rangle$, i.e.
  $\pi_{u,v}(gh^{-1}) \neq \pi_{u,v}(1_G)$. Right-translating by $h$
  gives $\pi_{u,v}(g) \neq \pi_{\rho(h)u, v}(h) = \pi_{u, v}(h)$ via
  the matrix-coefficient identity $\pi_{u,v}(gh) =
  \langle \rho(g)\rho(h)u, v\rangle = \pi_{\rho(h)u, v}(g)$.
\end{proof}

\section{Density of $\mathcal A$ via Stone--Weierstrass}

\begin{theorem}[$\mathcal A$ is dense in $C(G)$]
  \label{thm:c-mc-dense-CG}
  \lean{UnitaryRep.matrixCoeffSubalgebra_topologicalClosure_eq_top}
  % [NEW]
  \uses{cor:c-mc-sep-points, lem:c-mc-algebra}
  The closure of $\mathcal A$ in the uniform topology of $C(G,\mathbb C)$
  is all of $C(G, \mathbb C)$.
\end{theorem}

\begin{proof}
  \uses{cor:c-mc-sep-points}
  Apply the complex Stone--Weierstrass theorem
  \texttt{ContinuousMap.starSubalgebra\_topologicalClosure\_eq\_top\_of\_separatesPoints}
  (\texttt{Mathlib.Topology.ContinuousMap.StoneWeierstrass}) to
  $\mathcal A$. Hypotheses: $G$ compact (Definition~\ref{def:c-standing});
  $\mathcal A$ a star-subalgebra of $C(G, \mathbb C)$
  (Definition~\ref{def:c-mc-subalg}); $\mathcal A$ separates points
  (Corollary~\ref{cor:c-mc-sep-points}). Conclusion: $\mathcal A$ is
  dense in $C(G)$ in the uniform topology.
\end{proof}

\begin{corollary}[$\mathcal A$ is dense in $L^2(G)$]
  \label{cor:c-mc-dense-L2}
  \lean{UnitaryRep.matrixCoeffSubalgebra_topologicalClosure_eq_top_L2}
  % [~ML] Uniform density implies L^2 density on a finite measure space.
  \uses{thm:c-mc-dense-CG, thm:c-CG-dense, thm:c-Lp-inclusions}
  The closure of $\mathcal A$ in $L^2(G)$ (under the $L^2$ topology) is
  all of $L^2(G)$.
\end{corollary}

\begin{proof}
  \uses{thm:c-mc-dense-CG, thm:c-CG-dense, thm:c-Lp-inclusions}
  $\mathcal A$ is uniformly dense in $C(G)$ by
  Theorem~\ref{thm:c-mc-dense-CG}, hence $L^2$-dense in $C(G)$ since
  $\|\cdot\|_{L^2} \leq \|\cdot\|_{L^\infty}$ on a probability space
  (Theorem~\ref{thm:c-Lp-inclusions}). And $C(G)$ is $L^2$-dense in
  $L^2(G)$ (Theorem~\ref{thm:c-CG-dense}).
\end{proof}

\section{The isotypic decomposition of $L^2(G)$}

\begin{definition}[Isotypic component]
  \label{def:c-isotypic-comp}
  \lean{UnitaryRep.isotypicComponent}
  % [NEW]
  \uses{def:c-irrep, def:c-equiv-rep, def:c-regular-rep}
  Given an irreducible class $[\xi]$ (an isomorphism class of
  finite-dimensional continuous irreducible unitary representations of
  $G$), the $\xi$-\emph{isotypic component} of $L^2(G)$ (with respect to
  the right regular representation $\rho$) is
  $L^2(G)_\xi := \overline{\sum_{V \cong \xi} V}$,
  the closure of the sum of all closed right-$G$-invariant subspaces
  of $L^2(G)$ that are isomorphic to $\xi$.
\end{definition}

\begin{lemma}[Isotypic components are pairwise orthogonal]
  \label{lem:c-isotypic-orthog}
  \lean{UnitaryRep.isotypicComponent_orthogonal}
  % [NEW]
  \uses{def:c-isotypic-comp, cor:c-schur-distinct}
  $L^2(G)_\xi \perp L^2(G)_\eta$ whenever $[\xi] \neq [\eta]$.
\end{lemma}

\begin{proof}
  \uses{cor:c-schur-distinct}
  Pick $V \cong \xi$ and $W \cong \eta$ closed irreducible
  subrepresentations of $L^2(G)$. The orthogonal projection
  $P_W : L^2(G) \to W$ is bounded and $G$-equivariant
  (Lemma~\ref{lem:c-subrep-orthcomp}); restricting to $V$ gives an
  intertwiner $V \to W$, which is zero by Schur
  (Corollary~\ref{cor:c-schur-distinct}). Hence
  $V \perp W$. Since closures preserve orthogonality, the same holds
  for $L^2(G)_\xi$ and $L^2(G)_\eta$.
\end{proof}

\begin{theorem}[Peter--Weyl: $L^2$ decomposition]
  \label{thm:c-PW-L2-decomp}
  \lean{UnitaryRep.L2_isHilbertSum_isotypicComponents}
  % [NEW]
  \uses{cor:c-mc-dense-L2, lem:c-isotypic-orthog,
        cor:c-irr-finite-dim, def:c-isotypic-comp}
  $L^2(G)$ is the orthogonal Hilbert sum of the $\xi$-isotypic
  components, as $[\xi]$ ranges over all isomorphism classes of
  finite-dimensional continuous unitary irreducible representations of
  $G$:
  \[
     L^2(G) \;=\; \widehat{\bigoplus}_{[\xi]} L^2(G)_\xi.
  \]
  In particular every $f \in L^2(G)$ has a unique decomposition
  $f = \sum_{[\xi]} f_\xi$ with $f_\xi \in L^2(G)_\xi$, with
  $\sum \|f_\xi\|_2^2 = \|f\|_2^2$ (Parseval--Plancherel).
\end{theorem}

\begin{proof}
  \uses{cor:c-mc-dense-L2, lem:c-isotypic-orthog, def:c-isotypic-comp}
  By Lemma~\ref{lem:c-isotypic-orthog} the family
  $\{L^2(G)_\xi\}_{[\xi]}$ is an orthogonal family in $L^2(G)$. The
  closed sum $\overline{\sum_{[\xi]} L^2(G)_\xi}$ is a closed
  $G$-invariant subspace of $L^2(G)$. Each matrix coefficient
  $\pi_{u,v}$ of an irreducible $\xi$ lies in $L^2(G)_\xi$ (by the
  matrix-coefficient construction), so $\mathcal A$ is contained in the
  closed sum. By Corollary~\ref{cor:c-mc-dense-L2}, $\mathcal A$ is
  dense in $L^2(G)$, so the closed sum equals $L^2(G)$.

  In Lean: this is exactly the predicate
  \texttt{IsHilbertSum} from
  \texttt{Mathlib.Analysis.InnerProductSpace.l2Space}, applied to the
  orthogonal family of inclusion isometries
  $L^2(G)_\xi \hookrightarrow L^2(G)$.
\end{proof}

\begin{theorem}[Peter--Weyl: explicit form on each isotypic component]
  \label{thm:c-PW-each-isotypic}
  \lean{UnitaryRep.isotypicComponent_isHilbertSum_copies}
  % [NEW]
  \uses{thm:c-PW-L2-decomp, cor:c-irr-finite-dim, def:c-isotypic-comp}
  Each isotypic component $L^2(G)_\xi$ is itself an orthogonal Hilbert
  sum of $\dim V_\xi$-many copies of $V_\xi$ (so $L^2(G)_\xi$ has
  $\mathbb C$-dimension $(\dim V_\xi)^2$ when $\dim V_\xi < \infty$).
\end{theorem}

\begin{proof}
  \uses{thm:c-PW-L2-decomp, cor:c-irr-finite-dim}
  Standard. Use the (left, right) bi-equivariance of $L^2(G)$:
  $L^2(G)_\xi$ is also a closed left-$G$-invariant subspace, and as a
  representation under the joint $G \times G$-action it is isomorphic
  to $V_\xi \otimes V_\xi^*$. The matrix coefficient map
  $V_\xi \otimes V_\xi^* \to L^2(G)$, $u \otimes \ell \mapsto
  (g \mapsto \ell(\rho_\xi(g)^{-1} u))$, is an injective $G \times G$
  intertwiner with image $L^2(G)_\xi$; orthogonality
  $\langle u_1 \otimes \ell_1, u_2 \otimes \ell_2 \rangle =
  (\dim V_\xi)^{-1}\langle u_1, u_2 \rangle \overline{\langle \ell_1,
  \ell_2 \rangle}$ follows from Schur's orthogonality relations.
\end{proof}


\chapter{Schur Orthogonality, Plancherel, and Fourier Inversion}
\label{chap:c-orthogonality}

\section{Schur orthogonality of matrix coefficients}

\begin{theorem}[Schur orthogonality, distinct irreducibles]
  \label{thm:c-schur-orthog-distinct}
  \lean{UnitaryRep.matrixCoeff_orthog_of_not_equiv}
  % [NEW]
  \uses{cor:c-schur-distinct, def:c-matrix-coeff, lem:c-isotypic-orthog}
  Let $(\rho, V)$, $(\sigma, W)$ be finite-dimensional continuous unitary
  irreducibles with $\rho \not\cong \sigma$. Then for every
  $u_1, u_2 \in V$ and $w_1, w_2 \in W$,
  \[
     \int_G \langle \rho(g) u_1, u_2 \rangle_V \cdot
     \overline{\langle \sigma(g) w_1, w_2 \rangle_W} \, d\mu(g) = 0.
  \]
  Equivalently $\pi_{u_1, u_2}^\rho \perp \pi_{w_1, w_2}^\sigma$
  in $L^2(G)$.
\end{theorem}

\begin{proof}
  \uses{cor:c-schur-distinct, lem:c-isotypic-orthog}
  Direct from Lemma~\ref{lem:c-isotypic-orthog}: matrix coefficients of
  $\rho$ live in the $[\rho]$-isotypic component of $L^2(G)$, and those
  of $\sigma$ in the $[\sigma]$-component, which are orthogonal.
\end{proof}

\begin{theorem}[Schur orthogonality, same irreducible]
  \label{thm:c-schur-orthog-same}
  \lean{UnitaryRep.matrixCoeff_inner_eq}
  % [NEW]
  \uses{thm:c-schur-scalar, def:c-matrix-coeff}
  Let $(\rho, V)$ be a finite-dimensional continuous unitary irreducible
  with $d := \dim V$. Then for $u_1, u_2, v_1, v_2 \in V$,
  \[
     \int_G \langle \rho(g) u_1, v_1 \rangle \cdot
     \overline{\langle \rho(g) u_2, v_2 \rangle} \, d\mu(g)
     = \frac{1}{d}\,\langle u_1, u_2 \rangle \cdot
     \overline{\langle v_1, v_2 \rangle}.
  \]
\end{theorem}

\begin{proof}
  \uses{thm:c-schur-scalar}
  Standard.  Define a $G$-equivariant operator
  $T : V \to V$, $T u := \int_G \rho(g)\, u \otimes \overline{\rho(g)v_1} \cdot
  \pi_{u_2, v_2}(g)^* \, d\mu$ (compactness of $G$ ensures convergence).
  Use Theorem~\ref{thm:c-schur-scalar} to identify $T$ with a scalar,
  and compute the scalar by taking traces.
\end{proof}

\begin{corollary}[Orthonormal basis from matrix coefficients]
  \label{cor:c-mc-orthonormal-basis}
  \lean{UnitaryRep.matrixCoeff_hilbertBasis}
  % [NEW]
  \uses{thm:c-schur-orthog-same, thm:c-schur-orthog-distinct,
        thm:c-PW-L2-decomp}
  Let $\widehat{G}$ index a choice of one finite-dimensional unitary
  irreducible $V_\xi$ from each isomorphism class. For each $\xi$, fix
  an orthonormal basis $\{e^\xi_i\}_{i=1}^{d_\xi}$ of $V_\xi$. Then
  the family $\{\sqrt{d_\xi}\, \pi_{e^\xi_i, e^\xi_j}\}$ for
  $\xi \in \widehat G$, $1 \leq i, j \leq d_\xi$, is an
  orthonormal Hilbert basis of $L^2(G)$. In Lean: a
  \texttt{HilbertBasis ($\Sigma$ $\xi \in \widehat G$,
                        $\mathrm{Fin}\,d_\xi \times \mathrm{Fin}\,d_\xi$)
  ($L^2(G)$)}.
\end{corollary}

\begin{proof}
  \uses{thm:c-schur-orthog-same, thm:c-schur-orthog-distinct,
        thm:c-PW-L2-decomp}
  Combine Theorem~\ref{thm:c-PW-L2-decomp} (the matrix coefficients span
  a dense subspace) with
  Theorem~\ref{thm:c-schur-orthog-same} (giving the explicit
  $1/d_\xi$ normalisation factor) and
  Theorem~\ref{thm:c-schur-orthog-distinct} (orthogonality across
  $\xi$).
\end{proof}

\section{Fourier transform and Plancherel}

\begin{definition}[Operator-valued Fourier transform]
  \label{def:c-fourier}
  \lean{UnitaryRep.fourierTransform}
  % [NEW]
  \uses{def:c-matrix-coeff, def:c-isotypic-comp}
  For $f \in L^1(G)$ and an irreducible class $[\xi]$ with
  representative $(\rho_\xi, V_\xi)$, define
  $\widehat f(\xi) \in \mathrm{End}(V_\xi)$ by
  $\widehat f(\xi) := \int_G f(g)\, \rho_\xi(g) \, d\mu(g)$
  (Bochner integral in the finite-dimensional space
  $\mathrm{End}(V_\xi)$).
\end{definition}

\begin{theorem}[Plancherel for compact groups]
  \label{thm:c-plancherel}
  \lean{UnitaryRep.plancherel}
  % [NEW]
  \uses{thm:c-PW-each-isotypic, cor:c-mc-orthonormal-basis,
        def:c-fourier}
  For $f \in L^2(G)$:
  $\|f\|_{L^2(G)}^2 = \sum_{[\xi]} d_\xi \,\|\widehat f(\xi)\|_{HS}^2$,
  where $\|\cdot\|_{HS}$ is the Hilbert--Schmidt norm on
  $\mathrm{End}(V_\xi)$.
\end{theorem}

\begin{proof}
  \uses{cor:c-mc-orthonormal-basis, def:c-fourier}
  Expand $f$ in the orthonormal basis from
  Corollary~\ref{cor:c-mc-orthonormal-basis} and identify the resulting
  coefficients with the entries of $\widehat f(\xi)$.
\end{proof}

\begin{theorem}[Fourier inversion]
  \label{thm:c-fourier-inversion}
  \lean{UnitaryRep.fourierInversion}
  % [NEW]
  \uses{thm:c-plancherel, def:c-fourier}
  For $f \in L^2(G)$ (with $L^2$ convergence),
  $f(g) = \sum_{[\xi]} d_\xi \, \mathrm{tr}(\rho_\xi(g)^{-1} \widehat f(\xi))$.
\end{theorem}

\begin{proof}
  \uses{thm:c-plancherel, def:c-fourier}
  Same expansion as in Theorem~\ref{thm:c-plancherel}, summed back.
\end{proof}


\chapter{Full Statement of the Peter--Weyl Theorem}
\label{chap:c-statement}

The blueprint culminates in three equivalent forms of Peter--Weyl:

\begin{theorem}[Peter--Weyl, Form I: representations]
  \label{thm:c-PW-I}
  \lean{UnitaryRep.PeterWeyl_representations}
  % [NEW]
  \uses{thm:c-PW-L2-decomp, thm:c-PW-each-isotypic}
  Every continuous unitary representation of a compact group $G$ on a
  Hilbert space decomposes (uniquely) as a Hilbert sum of finite-dimensional
  irreducibles. In particular, all irreducible continuous unitary
  representations of $G$ are finite-dimensional.
\end{theorem}

\begin{proof}
  \uses{thm:c-PW-L2-decomp, thm:c-PW-each-isotypic, cor:c-irr-finite-dim}
  Apply the isotypic decomposition (Theorem~\ref{thm:c-PW-L2-decomp})
  componentwise. Finite-dimensionality of the irreducibles follows from
  Corollary~\ref{cor:c-irr-finite-dim} since every irreducible can be
  embedded in $L^2(G)$ by the matrix-coefficient construction (or
  directly by the regular representation).
\end{proof}

\begin{theorem}[Peter--Weyl, Form II: matrix coefficients]
  \label{thm:c-PW-II}
  \lean{UnitaryRep.PeterWeyl_matrixCoefficients_dense}
  % [NEW]
  \uses{thm:c-mc-dense-CG, cor:c-mc-orthonormal-basis}
  The set of matrix coefficients of finite-dimensional continuous unitary
  irreducible representations is uniformly dense in $C(G, \mathbb C)$,
  and the suitably normalised matrix coefficients form an orthonormal
  Hilbert basis of $L^2(G)$.
\end{theorem}

\begin{proof}\uses{thm:c-mc-dense-CG, cor:c-mc-orthonormal-basis}
  Combination of Theorem~\ref{thm:c-mc-dense-CG} (uniform density) and
  Corollary~\ref{cor:c-mc-orthonormal-basis} (Hilbert basis).
\end{proof}

\begin{theorem}[Peter--Weyl, Form III: $L^2$ decomposition]
  \label{thm:c-PW-III}
  \lean{UnitaryRep.PeterWeyl_L2_decomposition}
  % [NEW]
  \uses{thm:c-PW-L2-decomp, thm:c-PW-each-isotypic}
  As a representation of $G \times G$ under (left, right)-translation,
  \[
     L^2(G) \;\cong\;
     \widehat{\bigoplus}_{[\xi] \in \widehat G}
     V_\xi \boxtimes V_\xi^*,
  \]
  where each summand $V_\xi \boxtimes V_\xi^*$ is the
  finite-dimensional representation of $G \times G$ on
  $V_\xi \otimes V_\xi^*$ via $\rho_\xi \boxtimes \rho_\xi^*$.
\end{theorem}

\begin{proof}\uses{thm:c-PW-L2-decomp, thm:c-PW-each-isotypic}
  Combine Theorem~\ref{thm:c-PW-L2-decomp} and
  Theorem~\ref{thm:c-PW-each-isotypic}, the latter giving the
  identification of each isotypic component with
  $V_\xi \otimes V_\xi^*$ as a $G \times G$-representation.
\end{proof}


\chapter{Mathlib-Status Summary and Roadmap}
\label{chap:c-status}

\paragraph{What is already in Mathlib (\texttt{[ML]}, with \texttt{\textbackslash leanok}).}
The analytic foundations are mostly there:
\texttt{MeasureTheory.Measure.haarMeasure} and the \texttt{IsHaarMeasure}
typeclass with all invariance, regularity, and uniqueness results
(\texttt{Mathlib.MeasureTheory.Measure.Haar.\{Basic, Unique\}});
\texttt{MeasureTheory.Lp E 2 $\mu$} with its inner-product
structure (\texttt{Mathlib.MeasureTheory.Function.L2Space});
density of $C(G)$ in $L^p$ on a finite regular measure space
(\texttt{MeasureTheory.Lp.continuousMap\_dense});
the inclusions $L^p \subseteq L^q$ on a finite measure space; the bundled
\texttt{StarSubalgebra $\mathbb C$ C(G, $\mathbb C$)} type and the complex
Stone--Weierstrass theorem
(\texttt{ContinuousMap.starSubalgebra\_topologicalClosure\_eq\_top\_of\_separatesPoints});
\texttt{IsHilbertSum} and \texttt{HilbertBasis}
(\texttt{Mathlib.Analysis.InnerProductSpace.l2Space});
\texttt{IsCompactOperator} as a predicate
(\texttt{Mathlib.Analysis.Normed.Operator.Compact});
the multiplicative convolution \texttt{MeasureTheory.mconvolution}
(\texttt{Mathlib.Analysis.Convolution}).

\paragraph{Trivial wrappers (\texttt{[\textasciitilde ML]}).}
\texttt{UnitaryRep.matrixCoeffSubalgebra\_topologicalClosure\_eq\_top\_L2}
(uniform density implies $L^2$ density on a finite measure space).

\paragraph{Genuinely new contributions (\texttt{[NEW]}).}
There are roughly four blocks:

\begin{enumerate}

  \item \textbf{Continuous unitary representation theory.}
  \texttt{UnitaryRep $\mathbb C$ G H} (the bundled structure) and the
  basic API: \texttt{UnitaryRep.Subrep},
  \texttt{UnitaryRep.IsIrreducible}, \texttt{UnitaryRep.Intertwiner},
  \texttt{UnitaryRep.Equiv}, the orthogonal-complement closure of
  invariant subspaces (Lemma~\ref{lem:c-subrep-orthcomp}), and the
  Schur lemmas (Theorem~\ref{thm:c-schur-scalar},
  Corollary~\ref{cor:c-schur-distinct}). This is roughly the analog of
  what Mathlib's \texttt{Representation} infrastructure provides in the
  algebraic setting, and probably belongs in a new
  \texttt{Mathlib.RepresentationTheory.Unitary} file (or several).

  \item \textbf{Convolution as a compact self-adjoint operator on
  $L^2(G)$.} The chain
  Theorem~\ref{thm:c-conv-compact},
  Theorem~\ref{thm:c-conv-selfadjoint},
  Theorem~\ref{thm:c-conv-comm-right},
  Lemma~\ref{lem:c-approx-identity}.
  These are essentially folklore but are not currently in Mathlib for
  compact non-abelian groups.

  \item \textbf{The compact-operator spectral theorem
  (Theorem~\ref{thm:c-spectral}).}
  This is the most substantial gap.  Mathlib has the finite-dimensional
  version (\texttt{LinearMap.IsSymmetric.diagonalization}) but not the
  infinite-dimensional compact-self-adjoint case. Formalising it is a
  project in its own right (perhaps even larger than the rest of this
  blueprint combined). A workaround might be to prove just the weaker
  ``every compact self-adjoint operator has a non-zero eigenvalue with
  finite-dimensional eigenspace'', which is enough for
  Corollary~\ref{cor:c-irr-finite-dim} and
  Theorem~\ref{thm:c-irreducibles-separate}.

  \item \textbf{Peter--Weyl proper.} The decomposition theorems
  themselves: Theorem~\ref{thm:c-irreducibles-separate} (existence of
  separating finite-dim irreducibles), Theorem~\ref{thm:c-mc-dense-CG}
  (dense matrix coefficients via Stone--Weierstrass),
  Theorem~\ref{thm:c-PW-L2-decomp} ($L^2$ Hilbert sum),
  Theorem~\ref{thm:c-PW-each-isotypic} (per-isotypic structure),
  Schur orthogonality, Plancherel, Fourier inversion.

\end{enumerate}

\paragraph{Suggested formalisation order.}
\begin{enumerate}
  \item Define \texttt{UnitaryRep} and basic API
        (Chapter~\ref{chap:c-rep}).
  \item Convolution facts on a compact group
        (Chapter~\ref{chap:c-convolution}, sections 1).
  \item Either (a) formalise the compact-self-adjoint spectral theorem
        in full generality, or (b) extract the weaker ``every non-zero
        compact self-adjoint operator has a finite-dimensional
        eigenspace at a non-zero eigenvalue'' --- which is what we
        actually use --- as a focused mini-project.
  \item Schur's lemma (Chapter~\ref{chap:c-schur}).
  \item The matrix-coefficient subalgebra and Stone--Weierstrass
        (Chapter~\ref{chap:c-peter-weyl-L2}, sections 1--3).
  \item The Peter--Weyl decomposition and Plancherel
        (Chapter~\ref{chap:c-peter-weyl-L2}, sections 3--4 and
        Chapter~\ref{chap:c-orthogonality}).
\end{enumerate}

\paragraph{Comparison to the finite-group blueprint.}
Recall that for finite $G$ over an algebraically closed field $k$, the
proof reduces to ``Maschke + Wedderburn--Artin + algebraic closedness'',
with about \emph{four}--five genuinely new lemmas needed beyond Mathlib.
For a compact (possibly infinite) topological group, the analytic
infrastructure dominates and the new contributions are an order of
magnitude larger; the spectral theorem alone is a sub-project of
substantial size. This is in line with the historical observation that
Peter--Weyl for compact groups is much harder than the finite-group case
and required serious harmonic analysis to develop.
